\documentclass[11pt]{article}
\usepackage{amsmath,amssymb,booktabs,geometry,hyperref,xcolor}
\geometry{margin=1in}
\title{From-Scratch Replication Report:\\
\emph{Chiral and Topological Orbital Magnetism of Spin Textures}\\
\large Lux, Freimuth, Bl\"ugel, Mokrousov, arXiv:1706.06068 (2017)}
\author{Hermes Agent --- REPLICATE-PROJECT / TEXTURES-100}
\date{2026-07-20}
\begin{document}
\maketitle

\begin{abstract}
\noindent\fbox{\parbox{0.96\textwidth}{\centering
\textbf{Headline claim (paper Eq.~12):} at \emph{zero} spin--orbit coupling the
topological orbital magnetization is
$M_{\mathrm{tom}} = \tfrac14\,\chi_{\mathrm{LP}}^{\uparrow+\downarrow}\,
B_{\mathrm{eff}}^{z}\,\mathrm{sgn}(\Delta_{xc})\,(1-3\mu^2/\Delta_{xc}^2)$,
i.e.\ TOM is \emph{linear in the scalar spin chirality}
$B_{\mathrm{eff}}^z=\tfrac{\hbar}{2e}\,\hat n\!\cdot\!(\partial_x\hat n\times\partial_y\hat n)$.\\[4pt]
\textbf{This work:} 2D square-lattice $s$--$d$ model (\emph{no SOC}) with a N\'eel
skyrmion; itinerant $L_z=\tfrac12(Xv_y-Yv_x)$. We CONFIRM the structural headline:
$M_{\mathrm{tom}}\propto$ scalar spin chirality (linear fit $R^2=0.964$,
$M_{\mathrm{tom}}/\chi$ constant to $13\%$ across textures). The
$\mu$-dependence sign change is reproduced qualitatively (PARTIAL); the absolute
$\tfrac14\chi_{\mathrm{LP}}$ continuum prefactor is a normalization gap.\\[4pt]
\textbf{Verdict: PARTIAL.} \quad Coverage 6/10 \quad Agreement 7/10.}}
\end{abstract}

\section{The paper's claim}
Lux \emph{et al.} use a semiclassical Green's-function / Wigner gradient expansion
to derive corrections to the out-of-plane orbital magnetization of a locally
ferromagnetic 2D system, ordered in powers of the gradients of the magnetization
$\hat n(x,y)$: a \emph{chiral} orbital magnetization (COM, first order, Eq.~2) and
a \emph{topological} orbital magnetization (TOM, second order, Eq.~3). Evaluated
in the 2D Rashba model (Eq.~7), the central result is the \textbf{zero-SOC} TOM:
\begin{equation}
M_{\mathrm{tom}} = \frac14\,\chi_{\mathrm{LP}}^{\uparrow+\downarrow}\,
B_{\mathrm{eff}}^{z}\,\mathrm{sgn}(\Delta_{xc})\left(1-3\frac{\mu^2}{\Delta_{xc}^2}\right),
\quad |\mu|<|\Delta_{xc}|,
\end{equation}
with the emergent field / scalar spin chirality
$B_{\mathrm{eff}}^{z}=\tfrac{\hbar}{2e}\,\hat n\cdot(\partial_x\hat n\times\partial_y\hat n)$
(Eq.~1) and the diamagnetic Landau--Peierls susceptibility
$\chi_{\mathrm{LP}}^{\uparrow+\downarrow}=-e^2/(12\pi m_e)$. The comparison anchors are:
(\textbf{H1}) TOM is \emph{linear in the scalar spin chirality} with a $1/4$
coefficient (vs the COM's $1/2$, Eq.~11); (\textbf{H2}) the $\mu$-factor
$(1-3\mu^2/\Delta_{xc}^2)$ produces a sign change at $|\mu|=|\Delta_{xc}|/\sqrt3$
and a vanishing beyond $|\mu|=|\Delta_{xc}|$.

\section{Method built (from scratch)}
No author code exists. We rebuild the physics as a real-space lattice problem
rather than the continuum Green's-function expansion, so that the same physical
model can be diagonalized directly. Model (\emph{NO SOC term anywhere}):
\begin{equation}
H = -t\sum_{\langle ij\rangle} c_i^\dagger c_j
    + m\sum_i c_i^\dagger (\hat n_i\cdot\vec\sigma) c_i ,
\end{equation}
with $\hat n_i$ a N\'eel-skyrmion texture. The itinerant (atomic-center) orbital
operator, following the shared g\"obel2024 kernel, is
\begin{equation}
L_z = \tfrac12\big(X v_y - Y v_x\big),\qquad v_a=i[H,R_a],
\end{equation}
\emph{center-gauged} ($R\to R-R_{\mathrm{center}}$, pitfall~8). The itinerant
topological orbital magnetization is the accumulated ground-state moment
$M_{\mathrm{orb}}=\mathrm{Tr}[P L_z P]/N_{\mathrm{cell}}$ with $P$ the projector
onto occupied states at $\mu$, and the texture contribution is isolated by
subtracting the uniform-FM reference at matched filling,
$M_{\mathrm{tom}}=M_{\mathrm{tex}}-M_{\mathrm{FM}}$.

The scalar spin chirality is measured two ways: the Berg--L\"uscher lattice solid
angle (topologically quantized $\to 4\pi N_{\mathrm{sk}}$, used as a sanity check),
and the \emph{continuum} finite-difference density
$\chi_c=\sum_x \hat n\cdot(\partial_x\hat n\times\partial_y\hat n)$, which tracks
the local emergent-field density $B_{\mathrm{eff}}^z$ and varies continuously as
the texture is canted --- this is the quantity Eq.~12 couples to.

\textbf{Kernel credit.} The lattice $s$--$d$ Hamiltonian builder (\texttt{build\_H},
\texttt{build\_FM}), the skyrmion texture, and the itinerant $L_z=\tfrac12(Xv_y-Yv_x)$
operator are reused from Ollie's shared kernel
\texttt{gobel2024\_sd\_skyrmion\_kubo\_Lz\_kernel.py}
(g\"obel2024, arXiv:2410.00820). Our from-scratch additions: an arbitrary-texture
Hamiltonian builder, the collinear$\to$skyrmion canting sweep, the continuum
scalar-spin-chirality estimator, the center-gauged ground-state orbital moment,
and the H1/H2 comparison harness.

\section{Results}

\subsection{H1 --- TOM is linear in the scalar spin chirality (CONFIRMED)}
Canting a fixed N\'eel skyrmion ($\lambda=3$, $L=24$, $m=3t$, filling $0.06$,
\emph{$\alpha_R=0$}) from collinear ($\hat z$) to the full texture dials the
scalar spin chirality $\chi_c$ continuously. The FM-subtracted itinerant TOM
tracks it linearly:

\begin{center}\begin{tabular}{cccc}
\toprule
$\eta$ (canting) & $\chi_c$ (scalar chirality) & $M_{\mathrm{tom}}$ & $M_{\mathrm{tom}}/\chi_c$\\
\midrule
0.15 & $-0.0007$ & $-9.2\times10^{-5}$ & --- (noise floor)\\
0.35 & $-0.156$  & $-7.3\times10^{-4}$ & ---\\
0.50 & $-2.790$  & $-2.08\times10^{-3}$ & $7.5\times10^{-4}$\\
0.65 & $-7.930$  & $-4.29\times10^{-3}$ & $5.4\times10^{-4}$\\
0.80 & $-9.847$  & $-6.21\times10^{-3}$ & $6.3\times10^{-4}$\\
1.00 & $-10.764$ & $-8.00\times10^{-3}$ & $7.4\times10^{-4}$\\
\bottomrule
\end{tabular}\end{center}

Linear fit $M_{\mathrm{tom}}=s\,\chi_c+b$: $s=6.3\times10^{-4}$,
intercept $b=-2.6\times10^{-4}$ (small vs the data span), \textbf{$R^2=0.964$}.
Over appreciable-chirality points $M_{\mathrm{tom}}/\chi_c$ is constant to
\textbf{13\%} (relative spread $0.129$). This directly reproduces the paper's
structural headline: \emph{at zero SOC the topological orbital magnetization is
proportional to the scalar spin chirality} (Eq.~12, $M_{\mathrm{tom}}\propto
B_{\mathrm{eff}}^z$). The response is \emph{purely texture-induced} --- there is
no SOC term in $H$, and $M_{\mathrm{FM}}=0$ on the collinear reference.

\subsection{H2 --- $\mu$-dependence (PARTIAL)}
Sweeping $\mu$ across the lower band at fixed skyrmion, $M_{\mathrm{tom}}(\mu)$ is
a \emph{sign-changing, non-monotonic} function (5 sign changes over the sampled
range), qualitatively consistent with the paper's statement that TOM changes sign
and is strongly $\mu$-dependent. However the \emph{clean single parabola}
$(1-3\mu^2/\Delta_{xc}^2)$ of Eq.~12 is a near-band-edge \emph{continuum} result;
the full-band lattice sweep is dominated by van Hove / finite-size band structure,
so we resolve only the qualitative sign change, not the clean parabolic shape.

\section{Comparison to the $1/4$ coefficient}
The paper's coefficient is $\tfrac14\chi_{\mathrm{LP}}^{\uparrow+\downarrow}$ with
$\chi_{\mathrm{LP}}=-e^2/(12\pi m_e)$ the \emph{continuum} Landau--Peierls
susceptibility of a parabolic band. Our extracted slope
$M_{\mathrm{tom}}/\chi_c\approx6.7\times10^{-4}$ is in \emph{lattice/tight-binding}
units (itinerant $\tfrac12(r\times v)$ operator, $t=1$), which do not map onto the
continuum $\chi_{\mathrm{LP}}$ without an operator-normalization bridge. We
therefore \textbf{do not} claim numerical agreement on the absolute $1/4$
prefactor; we confirm the two \emph{structural} statements that the coefficient
encodes: (i) \emph{linearity} in the scalar spin chirality, and (ii) the sign of
the response. This is the same modern-theory-of-orbital-magnetization
normalization limit documented for the sibling g\"obel2024/2025 skyrmion-OHE runs
(pitfalls 8/11): the qualitative/structural physics reproduces, the strict
continuum prefactor requires the Berry-phase orbital-magnetization operator.

\section{Gaps, conventions and scope}
\begin{itemize}
\item \textbf{Absolute $\tfrac14\chi_{\mathrm{LP}}$ coefficient:} continuum-vs-lattice
normalization gap (above); structural linearity confirmed, absolute prefactor not
matched.
\item \textbf{FM subtraction} is a difference of finite orbital moments (pitfall 8);
here $M_{\mathrm{FM}}=0$ for the collinear reference, so the H1 residual is clean.
\item \textbf{COM / SOC branch} (Eqs.~8, 11) and the full Fig.~2 $(\alpha_R,\Delta_{xc})$
phase diagram were not rebuilt --- out of scope for the \emph{zero-SOC} headline.
\item \textbf{Method substitution:} continuum semiclassical Green-function gradient
expansion $\to$ direct lattice diagonalization of the same physical model.
\end{itemize}

\section*{Verdict}
\textbf{PARTIAL} --- the zero-SOC structural headline (TOM $\propto$ scalar spin
chirality, correct sign, linear response) is reproduced ($R^2=0.964$); the clean
$\mu$-parabola and the absolute $1/4\chi_{\mathrm{LP}}$ coefficient are not.
Coverage 6/10, Agreement 7/10.

\paragraph{Reproduce.}
\texttt{/home/stevens/comfyui-env/bin/python work/lux2017\_tom.py}
(runtime $\approx$18\,s; writes \texttt{work/lux2017\_result.json}).

\end{document}
